\uuid{vl2a}
\titre{Solution générale/raccord}
\auteur{hanani ; mzari}
\organisation{exo7}
\licence{CC BY-NC-SA 4.0}
\datecreate{2020-01-20}
\niveau{L1}
\module{Analyse}
\chapitre{Equation différentielle}
\sousChapitre{Résolution d'équation différentielle du premier ordre}
\difficulte{3}
\theme{Equations du premier ordre, solution générale, raccord}
\exo7qcmid{lille-2.equations-differentielles.018}

\contenu{
\begin{question}
\qclassification{225.01}
\qtags{motcle=solution générale/raccord}
On considère l'équation différentielle $(E)$ : $(x+1)y'+y=2x+1$. Parmi les affirmations suivantes, cocher celles qui sont vraies :
\begin{answers}
\good{La solution générale de $(E)$ sur $]-1,+\infty[$ est : $\displaystyle y=x+\frac{k}{x+1}$, $k\in \Rr$.}
\good{La solution générale de $(E)$ sur $]-\infty,-1[$ est :$\displaystyle y=x+\frac{k}{x+1}$, où $k\in \Rr$.}
\bad{Les solutions de $(E)$ sur $\Rr$ sont les fonctions $\displaystyle y=x+\frac{k}{x+1}$, où $k\in \Rr$.}
\bad{$(E)$ possède une infinité de solutions sur $\Rr$.}
\end{answers}
\begin{explanations}
Remarquer que $y_0=x$ est une solution particulière. Donc la solution générale de $(E)$ sur $]-1,+\infty[$, ou $]-\infty,-1[$, est : $\displaystyle y=x+\frac{k}{x+1}$, $k\in \Rr$. Si $y$ est une solution de $(E)$ sur $\Rr$, alors $y$ est une solution de $(E)$ sur $]-1,+\infty[$ et $y$ est une solution de $(E)$ sur $]-\infty,-1[$. De plus, avec $x=-1$ dans $(E)$, on aura $y(-1)=-1$. D'où
$$y(x)=\left\{\begin{array}{lll}\displaystyle y(x)=x+\frac{k_1}{x+1}&\mbox{si}&x\in ]-1,+\infty [ \\\\ \displaystyle y(x)=x+\frac{k_2}{x+1}&\mbox{si}&x\in ]-\infty ,-1[\\ -1&\mbox{si}&x=-1. \end{array}\right.$$
La continuité de $y$ en $-1$ implique que $k_1=k_2=0$. L'équation $(E)$ ne possède qu'une seule solution sur $\Rr$.
\end{explanations}
\end{question}
}
